%=====================================================================
\chapter{Text and Natural Language Processing}\label{ch:nlp}
%=====================================================================
\locator{4}{Part IV --- Deep Learning and Modern AI}

\begin{outcomes}
By the end of this chapter you will be able to:
\begin{tight}
  \item \textbf{Apply} the standard text preprocessing pipeline and justify each step.
  \item \textbf{Compute} TF-IDF and explain what it measures.
  \item \textbf{Distinguish} sparse bag-of-words representations from dense embeddings.
  \item \textbf{Build} a text classifier and evaluate it appropriately.
  \item \textbf{Interpret} topic model output and state its limitations.
\end{tight}
\end{outcomes}

\begin{prereq}
\chref{math} (cosine similarity --- the standard text metric),
\chref{classification} (naive Bayes) and \chref{architectures} (embeddings and
attention).
\end{prereq}

\begin{keyterms}
corpus $\bullet$ token $\bullet$ tokenisation $\bullet$ stop words $\bullet$
stemming $\bullet$ lemmatisation $\bullet$ bag of words $\bullet$ n-gram $\bullet$
term frequency $\bullet$ inverse document frequency $\bullet$ TF-IDF $\bullet$
sparse vs dense $\bullet$ word embedding $\bullet$ contextual embedding $\bullet$
topic modelling $\bullet$ named entity recognition
\end{keyterms}

\section{Turning Words into Numbers}

Every model in this book consumes an $n \times p$ matrix of numbers
(\dsref{matrix-anatomy}). Text is not that. The entire first half of natural
language processing is the problem of producing that matrix without throwing away
the meaning.

\dsfig{%
\begin{tikzpicture}[font=\scriptsize,
  st/.style={rounded corners=3pt, draw=#1, fill=#1!9, text=#1!50!black,
             text width=2.15cm, align=center, font=\tiny\bfseries, minimum height=1.0cm},
  arr/.style={-{Stealth[length=2mm]}, gray!60, line width=0.9pt},
  ex/.style={font=\tiny\ttfamily, text=gray!50!black, align=center, text width=2.4cm}]

\node[st=secondary] (a) at (0,0)     {1. RAW TEXT};
\node[st=primary]   (b) at (2.65,0)  {2. NORMALISE\\{\tiny lowercase, strip}};
\node[st=tealhl]    (c) at (5.3,0)   {3. TOKENISE\\{\tiny split into units}};
\node[st=goldhl]    (d) at (7.95,0)  {4. FILTER\\{\tiny stop words}};
\node[st=purplehl]  (e) at (10.6,0)  {5. STEM /\\LEMMATISE};
\node[st=accent]    (f) at (13.25,0) {6. VECTORISE};

\foreach \x/\y in {a/b,b/c,c/d,d/e,e/f}{\draw[arr] (\x) -- (\y);}

\node[ex] at (0,-1.35)     {"The Servers\\CRASHED again!"};
\node[ex] at (2.65,-1.35)  {the servers\\crashed again};
\node[ex] at (5.3,-1.35)   {[the, servers,\\crashed, again]};
\node[ex] at (7.95,-1.35)  {[servers,\\crashed, again]};
\node[ex] at (10.6,-1.35)  {[server,\\crash, again]};
\node[ex] at (13.25,-1.35) {[0, 1, 0, 2,\\1, 0, ...]};

\node[anchor=west, align=left, text width=13.4cm, font=\tiny, text=accent!75!black]
     at (-1.1,-2.55)
     {\textbf{Every step discards information, so every step is a decision.} Removing stop
      words destroys negation --- ``not vulnerable'' becomes ``vulnerable''. Lowercasing
      merges \texttt{US} and \texttt{us}, and in security merges an IP-adjacent token with
      a common word. Stemming maps \emph{university} and \emph{universal} to the same
      stem. Apply each step only if you can say what it buys you.};
\end{tikzpicture}}%
{The classical text preprocessing pipeline. Modern transformer models skip most of
steps 2--5, using subword tokenisation directly --- but the classical pipeline is still
what you want for fast, explainable, low-resource classifiers.}{nlp-pipeline}

\section{Bag of Words and TF-IDF}

\begin{definitionbox}[Bag of Words]
Represent each document as a vector of word counts over a fixed vocabulary. Word
order is discarded entirely --- hence ``bag''. Crude, and surprisingly effective for
classification.
\end{definitionbox}

\begin{definitionbox}[TF-IDF]
\[
\text{tf-idf}(t, d) \;=\; \underbrace{\text{tf}(t,d)}_{\text{count in this doc}}
\times
\underbrace{\log\frac{N}{\text{df}(t)}}_{\text{rarity across the corpus}}
\]
A term scores highly when it is frequent \emph{in this document} and rare
\emph{across the corpus} --- which is a good working definition of ``distinctive''.
\end{definitionbox}

\begin{worked}[Computing TF-IDF]
A corpus of $N = 1000$ support tickets. Consider one ticket of 100 words.

\smallskip
\begin{center}
\begin{tabular}{@{}lrrrr@{}}
\toprule
\textbf{Term} & \textbf{count} & \textbf{tf} & \textbf{docs containing} & \textbf{idf} \\
\midrule
the         & 8 & 0.08 & 1000 & $\log(1000/1000)=0$ \\
server      & 4 & 0.04 & 300  & $\log(1000/300)=1.20$ \\
kerberos    & 2 & 0.02 & 12   & $\log(1000/12)=4.42$ \\
\bottomrule
\end{tabular}
\end{center}

\smallskip
\textbf{TF-IDF scores:}
\begin{tight}
  \item \texttt{the}: $0.08 \times 0 = \mathbf{0}$ --- appears everywhere, so it
        distinguishes nothing. Note that IDF has removed it \emph{automatically},
        without any stop-word list.
  \item \texttt{server}: $0.04 \times 1.20 = \mathbf{0.048}$
  \item \texttt{kerberos}: $0.02 \times 4.42 = \mathbf{0.088}$
\end{tight}

\smallskip
\textbf{The result that matters:} \texttt{kerberos} outranks \texttt{server}
despite appearing half as often, because it is rare across the corpus and therefore
tells you far more about which ticket this is. That is the whole idea, and it is
why TF-IDF remains a strong baseline forty years on.
\end{worked}

\section{From Sparse to Dense}

\begin{alertbox}[The Weakness of Bag of Words]
In a TF-IDF space, \emph{server}, \emph{host} and \emph{machine} are three
completely unrelated dimensions --- as unrelated as \emph{server} and
\emph{banana}. The representation has no notion of meaning, only of spelling. A
model trained on documents saying ``server'' will learn nothing applicable to
documents saying ``host''.
\end{alertbox}

\begin{definitionbox}[Word Embedding]
A dense vector of a few hundred dimensions, learned so that words appearing in
similar contexts get similar vectors. Meaning becomes geometry: cosine similarity
between embeddings measures semantic relatedness.
\end{definitionbox}

\dsfig{%
\begin{tikzpicture}
\begin{groupplot}[
  group style={group size=2 by 1, horizontal sep=1.6cm},
  width=6.4cm, height=4.6cm,
  axis lines=left, tick label style={font=\tiny}, label style={font=\tiny},
  title style={font=\scriptsize\bfseries},
]
\nextgroupplot[title={\color{accent}Sparse (TF-IDF): 20{,}000 dimensions},
  xlabel={vocabulary index}, ylabel={weight},
  xmin=0, xmax=40, ymin=0, ymax=1.1, ybar, bar width=2.2pt, xtick=\empty]
\addplot[draw=accent, fill=accent!35] coordinates {
 (4,0.62)(11,0.31)(19,0.88)(28,0.44)(35,0.21)};
\node[font=\tiny, text=gray!55!black, align=center, anchor=north] at (axis cs:20,1.05)
     {mostly zeros; each dimension\\is one exact word};

\nextgroupplot[title={\color{greenhl}Dense (embedding): 300 dimensions},
  xlabel={dimension}, ylabel={value},
  xmin=0, xmax=40, ymin=-1.0, ymax=1.15, ybar, bar width=2.2pt, xtick=\empty]
\addplot[draw=greenhl, fill=greenhl!35] coordinates {
 (1,0.42)(2,-0.31)(3,0.18)(4,0.66)(5,-0.52)(6,0.24)(7,-0.14)(8,0.71)(9,-0.38)
 (10,0.09)(11,0.55)(12,-0.62)(13,0.33)(14,0.12)(15,-0.45)(16,0.28)(17,0.61)
 (18,-0.22)(19,0.37)(20,-0.51)(21,0.44)(22,0.16)(23,-0.33)(24,0.58)(25,-0.19)
 (26,0.26)(27,0.48)(28,-0.41)(29,0.13)(30,0.35)(31,-0.28)(32,0.52)(33,0.07)
 (34,-0.46)(35,0.39)(36,0.21)(37,-0.35)(38,0.63)(39,-0.11)(40,0.29)};
\node[font=\tiny, text=gray!55!black, align=center, anchor=north] at (axis cs:20,1.1)
     {every dimension used; similar words\\get similar vectors};
\end{groupplot}
\end{tikzpicture}}%
{Sparse versus dense text representations. TF-IDF is interpretable but has no concept of
synonymy; embeddings capture meaning but no single dimension means anything you can
name.}{sparse-dense}

\begin{conceptbox}[Static versus Contextual Embeddings]
Early embeddings (word2vec, GloVe) give each word one fixed vector, so
\emph{bank} in ``river bank'' and ``bank account'' share a vector --- a real
limitation. Transformer models (\chref{architectures}) produce \emph{contextual}
embeddings: the vector for a word depends on the sentence it is in. This is the
main practical advance behind modern NLP, and it is what \chref{genai} builds on.
\end{conceptbox}

\section{Topic Modelling}

\begin{definitionbox}[Topic Modelling]
An unsupervised method (\chref{unsupervised} applied to text) that discovers
recurring themes in a corpus. Each topic is a distribution over words; each
document is a mixture of topics. Latent Dirichlet Allocation (LDA) is the classical
method.
\end{definitionbox}

\dsfig{%
\begin{tikzpicture}[font=\scriptsize,
  t/.style={rounded corners=3pt, draw=#1, fill=#1!8, text width=3.2cm, align=left,
            font=\tiny, inner sep=5pt},
  hd/.style={rounded corners=3pt, fill=#1, text=white, font=\tiny\bfseries,
             text width=3.2cm, align=center, minimum height=0.55cm}]

\node[hd=secondary] at (0,0)   {Topic 1 (auto-named ``login'')};
\node[t=secondary]  at (0,-1.3){%
  \textbf{0.081}~password\\ \textbf{0.064}~reset\\ \textbf{0.052}~locked\\
  \textbf{0.041}~account\\ \textbf{0.033}~expired};

\node[hd=greenhl] at (3.7,0)   {Topic 2 (``network'')};
\node[t=greenhl]  at (3.7,-1.3){%
  \textbf{0.077}~vpn\\ \textbf{0.069}~connection\\ \textbf{0.048}~timeout\\
  \textbf{0.040}~dropped\\ \textbf{0.031}~wifi};

\node[hd=goldhl] at (7.4,0)    {Topic 3 (``hardware'')};
\node[t=goldhl]  at (7.4,-1.3) {%
  \textbf{0.073}~laptop\\ \textbf{0.061}~screen\\ \textbf{0.050}~battery\\
  \textbf{0.038}~replace\\ \textbf{0.029}~keyboard};

\node[hd=purplehl] at (11.1,0)  {Topic 4 (``???'')};
\node[t=purplehl]  at (11.1,-1.3){%
  \textbf{0.045}~please\\ \textbf{0.039}~thanks\\ \textbf{0.036}~urgent\\
  \textbf{0.031}~asap\\ \textbf{0.028}~regards};

\node[anchor=west, align=left, text width=13.4cm, font=\tiny, text=accent!75!black]
     at (-1.75,-3.1)
     {\textbf{Topic 4 is the lesson.} LDA found a genuine statistical pattern --- politeness
      formulae --- that is not a topic in any useful sense. Topic models produce
      \emph{word clusters}, and it is entirely on you to decide which are meaningful. The
      labels above were written by a human; the algorithm supplied only the numbers.};
\end{tikzpicture}}%
{Output of a four-topic LDA on IT support tickets. Three topics are useful; the fourth is
a statistical artefact. Never present topic model output without human
inspection.}{topic-model}

%---------------------------------------------------------------------
\begin{fourlenses}{Text Analytics}
\lensLA{Analyse free-text reflections, forum posts and module feedback. Sentiment
and topic trends across a semester reveal problems that numeric engagement data
misses entirely --- a cohort can be clicking dutifully and telling you in prose that
they are lost. \textbf{Two cautions:} student writing is identifiable even after
anonymisation, and automated sentiment scoring of individuals is an ethical
decision, not a technical one (\chref{ethics}).}

\lensPM{Classify support tickets and requirement documents automatically, and mine
retrospective notes for recurring themes. TF-IDF plus logistic regression is
usually sufficient and takes an afternoon; it is also explainable, which matters
when the classifier is routing work to people. Topic modelling over several years of
retrospectives is a genuinely useful way to find problems the organisation keeps
re-discovering.}

\lensIOT{Text appears where you might not expect it: device logs, error strings and
maintenance notes. Treating log lines as text --- tokenising, then clustering
templates --- turns an unmanageable stream into a manageable set of message types,
and a \emph{new} template appearing is itself an anomaly signal. Engineers'
free-text maintenance records are often the only ground-truth labels available for
predictive maintenance (\chref{iot}).}

\lensSEC{Phishing and spam detection is the classic application: naive Bayes on
TF-IDF remains fast enough for mail-gateway volumes and is still competitive.
Beyond that, text classification applies to malware strings, command-line arguments
and DNS names. Note the adversarial character: spammers deliberately misspell,
insert homoglyphs and pad with innocuous text specifically to defeat tokenisation
--- so preprocessing choices here are a security control, not a convenience.}
\end{fourlenses}

%---------------------------------------------------------------------
\begin{lab}[TF-IDF Classifier and Topics]
\begin{lstlisting}[language=Python]
from sklearn.feature_extraction.text import TfidfVectorizer
from sklearn.linear_model import LogisticRegression
from sklearn.pipeline import make_pipeline
from sklearn.decomposition import LatentDirichletAllocation
from sklearn.metrics import classification_report
import numpy as np

# --- 1. a strong, explainable text classifier in four lines -----------
clf = make_pipeline(
    TfidfVectorizer(ngram_range=(1, 2),      # unigrams + bigrams: "not vulnerable"
                    min_df=3,                # ignore terms in fewer than 3 docs
                    max_df=0.85,             # ignore terms in >85% of docs
                    sublinear_tf=True),
    LogisticRegression(max_iter=1000, class_weight="balanced"),
).fit(train_texts, train_labels)

print(classification_report(test_labels, clf.predict(test_texts), digits=3))

# --- 2. WHY did it decide that? the great virtue of this approach -----
vec, lr = clf.named_steps["tfidfvectorizer"], clf.named_steps["logisticregression"]
terms = np.array(vec.get_feature_names_out())
top = lr.coef_[0].argsort()
print("most indicative of the positive class:", terms[top[-12:]][::-1])
print("most indicative of the negative class:", terms[top[:12]])

# --- 3. topic modelling: ALWAYS read the topics before believing them -
counts = TfidfVectorizer(max_df=0.85, min_df=5).fit(all_texts)
lda = LatentDirichletAllocation(n_components=6, random_state=0)
lda.fit(counts.transform(all_texts))
vocab = counts.get_feature_names_out()
for k, comp in enumerate(lda.components_):
    print(f"topic {k}:", ", ".join(vocab[i] for i in comp.argsort()[-8:][::-1]))
    # Inspect each one. Expect at least one to be an artefact like Topic 4 above.
\end{lstlisting}
\end{lab}

%---------------------------------------------------------------------
\begin{checkpoint}
\begin{tightnum}
  \item A word appears in every document of a corpus. What is its IDF, and what
        does that do to its TF-IDF score?
  \item Give a concrete example where removing stop words would reverse the meaning
        of a document.
  \item Why does TF-IDF treat ``server'' and ``host'' as completely unrelated, and
        what representation fixes it?
\end{tightnum}
\end{checkpoint}

%---------------------------------------------------------------------
\begin{chaptersummary}
\begin{tight}
  \item Text must become a numeric matrix; every preprocessing step trades
        information for tractability and must be justified.
  \item Bag of words discards order; TF-IDF weights terms by frequency in the
        document and rarity in the corpus, and downweights common words
        automatically.
  \item Sparse representations are interpretable but blind to synonymy; dense
        embeddings capture meaning as geometry but individual dimensions mean
        nothing.
  \item Contextual embeddings from transformers give a word a different vector in
        each sentence, which is the main modern advance.
  \item Topic models find word clusters, not topics. Human inspection is mandatory.
  \item TF-IDF plus logistic regression remains a fast, strong, explainable baseline
        --- try it before anything larger.
\end{tight}
\end{chaptersummary}

\begin{reviewq}
  \item \textbf{(MCQ)} A term appearing in 1 of 1000 documents has an IDF of roughly:
        (a)~0 (b)~1 (c)~3 (d)~7.
  \item \textbf{(MCQ)} Which representation captures that ``laptop'' and ``notebook''
        are related? (a)~bag of words (b)~TF-IDF (c)~word embeddings (d)~n-grams.
  \item \textbf{(Short)} Explain in your own words what TF-IDF measures and why the
        IDF term makes stop-word removal partly redundant.
  \item \textbf{(Short)} Give two reasons why bigrams often improve a text classifier.
  \item \textbf{(Short)} Why must topic model output be inspected by a human before
        it is reported?
  \item \textbf{(Applied)} Design a pipeline to classify 40{,}000 IT support tickets
        into eight categories. State your preprocessing choices with a justification
        for each, your representation, your model, and your evaluation metric ---
        noting that the categories are heavily imbalanced.
  \item \textbf{(Applied)} A phishing classifier trained on TF-IDF starts failing.
        Investigation shows attackers now write ``pa55word'' and ``inv0ice''. Explain
        why the model failed, and propose two changes that would make it more robust.
\end{reviewq}
